\chapter{Exigences et fonctionnalités}
\label{ch:exigences}

\section{Exigences pédagogiques (rappel)}

Le projet vise à satisfaire les critères suivants, tels que formulés dans l'énoncé du cours :
\begin{itemize}[leftmargin=*]
  \item \textbf{Blazor Web Application} : utilisation de Blazor en mode interactif serveur (\texttt{InteractiveServer}) pour les pages nécessitant des interactions riches (formulaires, graphiques, thème).
  \item \textbf{Architecture propre} : séparation des responsabilités entre UI, services métier et persistance.
  \item \textbf{Entity Framework Core} : approche \emph{Code First} avec \textbf{migrations} versionnées.
  \item \textbf{CRUD} complet sur les entités principales (actifs, transactions ; budget ajustable).
  \item \textbf{Recherche et filtrage} dynamiques via LINQ (actifs : texte, type ; transactions : pagination et tri).
  \item \textbf{Tableau de bord analytique} : indicateurs, graphiques (répartition, courbe de valeur dans le temps).
  \item \textbf{Validation} : \emph{Data Annotations} sur les entités et formulaires Blazor.
  \item \textbf{Asynchronisme} : méthodes \texttt{async}/\texttt{await} côté services et accès EF (\texttt{ToListAsync}, \texttt{SaveChangesAsync}, etc.).
\end{itemize}

\section{Fonctionnalités métier couvertes}

\begin{table}[h]
  \centering
  \caption{Fonctionnalités par domaine}
  \begin{tabular}{@{}p{4.2cm}p{10cm}@{}}
    \toprule
    \textbf{Domaine} & \textbf{Description} \\
    \midrule
    Actifs & CRUD, symbole unique, type Stock/Crypto, prix courant ; liste paginée, tri, recherche avec temporisation (\emph{debounce}). \\
    Transactions & Saisie achat/vente, calcul du montant total, contrôle de liquidité et de quantité détenue ; historique paginé ; export CSV. \\
    Budget & Consultation et mise à jour du budget initial et de la liquidité disponible. \\
    Portefeuille & Instantané : valeur titres, cash, valeur totale, gain/perte vs budget initial, ROI, répartition par actif et par type. \\
    Historique valeur & Points journaliers (\texttt{PortfolioValuePoint}) mis à jour à la visite du tableau de bord et après transactions. \\
    Marchés (optionnel) & Mode \emph{simulé}, mode \emph{API} (Finnhub + CoinGecko avec cache), import CSV des prix. \\
    Alertes & Seuils configurables : perte vs budget, liquidité basse, volatilité élevée par actif. \\
    Data science & SMA, volatilité annualisée approximative, Sharpe-like, drawdown, signal tendance, backtest SMA. \\
    UX & Thème clair/sombre (Bootstrap), graphiques Chart.js, navigation par menu. \\
    \bottomrule
  \end{tabular}
\end{table}

\section{Extension « créativité »}

L'extension retenue apporte une \textbf{valeur analytique et décisionnelle} complémentaire :
\begin{itemize}[leftmargin=*]
  \item quantification du \textbf{risque} (volatilité, drawdown) ;
  \item \textbf{tendance} par rapport à une moyenne mobile ;
  \item \textbf{comparaison risque/rendement} via un indicateur de type Sharpe simplifié ;
  \item \textbf{simulation historique} d'une règle d'investissement (backtest) pour argumenter à l'oral.
\end{itemize}
Cette extension reste cohérente avec le domaine « portefeuille » et s'appuie sur un historique de prix (généré si absent lors du premier chargement analytique).
